\section{Conclusion}
\label{sec:conclusion}

\subsection{Summary}
In this research, we have developed a novel data augmentation framework, designated as the Conditional Auxiliary Classifier Cycle-consistent Generative Adversarial Network with Wasserstein Gradient Penalty (CAC-CycleGAN-WGP), to address the persistent challenges of imbalanced bearing fault diagnosis. The proposed architecture integrates a conditional auxiliary classifier to enhance class-conditional generation, while the inclusion of the Wasserstein distance with a gradient penalty (WGP) stabilizes the training mechanism and effectively mitigates the risks of gradient vanishing. Furthermore, we have introduced a Stacked Autoencoder (SAE)-based evaluation methodology to quantitatively assess the fidelity and diversity of the generated vibration samples. By mapping high-dimensional vibration signals into a compact and representative latent feature space, the SAE facilitates a more robust feature-level comparison between synthetic and real experimental data. This ensures that the generated minority class samples accurately represent the underlying physical fault characteristics, providing a solid foundation for robust diagnostic model training.

\subsection{Key Findings}
Comprehensive experimental evaluations conducted on various standard bearing fault datasets demonstrate the exceptional effectiveness of the CAC-CycleGAN-WGP framework. Our results indicate that the proposed methodology significantly alleviates the performance degradation typically caused by extreme data imbalance scenarios in industrial settings. Specifically, even when subjected to a challenging imbalance ratio of 1:20, the framework maintains an impressively high level of diagnostic accuracy, consistently achieving an F-1 score exceeding 0.85. This performance markedly outperforms five state-of-the-art baseline models, including traditional GAN configurations and various SMOTE-based variants. A critical finding is the framework's superior ability to prevent mode collapse through the synergy of the cycle-consistent loss function and the auxiliary classifier constraints. Detailed ablation studies further underscore the vital contribution of the CAC module, which independently accounts for an approximately 40\% improvement in minority class classification accuracy. Additionally, the developed model exhibits remarkable noise robustness, maintaining high and stable diagnostic performance across a wide spectrum of Signal-to-Noise Ratios (SNR) ranging from -5~dB to 20~dB. Furthermore, across various experimental tests involving different rotational speeds and load conditions, the diagnostic precision variance remains consistently below 2.3\%, underscoring the strong generalizability of the proposed approach.

\subsection{Future Work}
While the CAC-CycleGAN-WGP framework has shown highly promising results in controlled environments, several critical avenues for future research remain. Subsequent efforts will focus on extending these augmentation capabilities to more complex non-stationary operating conditions where rotational speeds are time-varying, posing greater challenges for feature extraction and signal generation. Furthermore, we aim to investigate the framework's practical deployment in complex and harsh industrial settings characterized by multi-source interference and varying sensor locations. Integrating advanced domain adaptation techniques into the generative process could further enhance the transferability of the augmented data across different machines and sensor configurations. Finally, optimizing the computational efficiency of the training process to allow for real-time diagnostic model updates in edge computing devices remains a high priority for ensuring the practical, large-scale applicability of this technology within the evolving smart manufacturing and predictive maintenance ecosystems.
